\chapter{Results and Analysis}
\label{chap:results}

This chapter presents the complete experimental results with mechanistic analysis of why each contribution improves performance and what general lessons can be extracted. Results are grouped by theme: Section~\ref{sec:dlvsml} analyzes why deep learning outperforms machine learning baselines; Section~\ref{sec:feature_results} evaluates the multi-bandwidth and ANM channel contributions; Section~\ref{sec:struct_results} examines the structural attention bias; Section~\ref{sec:2d_results} presents the 2D pipeline ablations; Section~\ref{sec:full_ablation} provides the cumulative ablation; and Section~\ref{sec:lessons} synthesizes the general principles that emerged from this work.

% =====================================================================
\section{Deep Learning vs Machine Learning}
\label{sec:dlvsml}
% =====================================================================

This section establishes the mechanistic reasons why deep learning outperforms machine learning on this task, with implications for how future causal discovery systems should be designed.

\subsection{v10: LightGBM on Scalar Features}

Version v10 trains a LightGBM classifier on eleven hand-crafted scalar features including Pearson correlation, Spearman correlation, mutual information, and ANM scores, achieving 63.11\% Balanced Accuracy. This result is substantially below every Conv1D-based model in the thesis. The fundamental limitation is not the number of features or the choice of classifier but the irreversible nature of scalar compression: every feature in v10 reduces $n = 1000$ observations to a single number, discarding the functional form of the conditional relationship between variables.

Consider Mediator ($X \to K \to Y$) versus Consequence of X ($X \to K$ only). Both roles have the edge $X \to K$ in their DAG structure, creating dependence between $X$ and $K$ through the same mechanism. The distinguishing signal is whether $K$ subsequently affects $Y$, which is reflected in the curvature, heteroscedasticity, and tail behavior of the regression curve $\hat{m}_{X \to K}$ and the chain regression $\hat{m}_{K \to Y}$. No scalar summary of these curves -- whether correlation, mutual information, or ANM score -- can capture this discriminating information. The scalar encodes only the strength of the dependence, not its functional form.

\subsection{v13: Full ML Fullstack Pipeline}

Version v13 represents the maximum reasonable engineering effort under the ML paradigm: 300+ engineered features including correlation coefficients, partial correlations, conditional mutual information, HSIC, ANM scores, and outputs from four classical causal discovery algorithms (PC, LiNGAM, NOTEARS, GES), processed through an XGBoost/LightGBM/CatBoost ensemble with GNN refinement. This achieves 72.64\% -- below the 3-channel DL reimplementation baseline (v2: 73.96\%) and 8.44 percentage points below the proposed model.

The gap is mechanistic, not a matter of insufficient feature engineering. Every scalar feature in v13 is a fixed aggregation of the $n = 1000$ observations, and the information loss from this compression is irreversible. The spatial structure of the joint density, the shape of the conditional expectation function, the presence of multimodality or heteroscedasticity -- all the information that distinguishes the eight causal roles -- is discarded when observations are reduced to scalars. In contrast, the 8-channel edge tensor (channels 2--8) retains the full length-$n$ curve, and Conv1D learns which aspects of that curve are discriminative. The functional form of the causal relationship, not its aggregate strength, is the discriminating signal, and it is the information that scalar compression irrevocably destroys.

Table~\ref{tab:dl_vs_ml} summarizes the comparison. Machine learning approaches have genuine advantages -- interpretability, no GPU requirement, faster development cycles -- that make them appropriate for many tasks. For this specific problem, however, the discriminating signal lives in the shape of nonlinear functional relationships rather than in aggregate statistics, making the information loss from scalar compression decisive.

\begin{table}[ht]
	\centering
	\caption{Deep learning vs.\ machine learning on the ADIA Lab Challenge. Scalar compression (v13: 72.64\%) is 8.55 points below the proposed model (v26e+xyaug: 81.19\%) because the functional form of conditional relationships -- the discriminating signal for the eight causal roles -- is destroyed when $n=1000$ observations are reduced to scalar aggregates.}
	\label{tab:dl_vs_ml}
	\small
	\begin{tabularx}{\textwidth}{lXcc}
		\toprule
		\textbf{Approach} & \textbf{Description} & \textbf{BA} & \textbf{Gap} \\
		\midrule
		ML minimal (v10) & 11 scalar features & 63.11\% & $-18.08$ \\
		ML fullstack (v13) & 300+ scalars + algo outputs & 72.64\% & $-8.55$ \\
		DL proposed (v26e+xyaug) & 8-ch, $n=1000$ curves & \textbf{81.19\%} & --- \\
		\bottomrule
	\end{tabularx}
\end{table}

% =====================================================================
\section{Feature Engineering Contributions}
\label{sec:feature_results}
% =====================================================================

This section analyzes the incremental contributions of each channel in the 8-channel edge tensor, establishing that multi-bandwidth kernel regression and ANM residuals each provide distinct and complementary improvements.

\subsection{Multi-Bandwidth Kernel Regression Channels ($+1.48\%$)}

Extending the 3-channel baseline (v2: 73.96\%) by adding multi-bandwidth kernel regression channels at $h \in \{0.2, 1.0\}$ alongside the original $h = 0.5$ produces v5m at 75.44\%, a gain of $+1.48\%$. Each bandwidth captures a different smoothness scale of the nonlinear conditional expectation: $h = 0.2$ is sensitive to sharp local transitions and fine nonlinear structure, while $h = 1.0$ reflects global monotone trends and is robust to local noise. Using all three bandwidths together prevents the model from depending on any single fixed choice of $h$, which is critical because the optimal smoothness varies across different causal mechanisms and data generating processes. The gain is confirmed on the public leaderboard: v5m+xyaug achieves 74.80\% LB, a $+1.97\%$ improvement over the corresponding baseline, confirming that the improvement reflects genuine generalization rather than overfitting to the local test split.

The multi-bandwidth channels address a fundamental limitation of single-bandwidth representations: no fixed bandwidth is optimal for all causal relationships. Some variable pairs have sharp threshold-like transitions (requiring small $h$), while others have smooth monotone relationships (better captured by large $h$). The ANM residual channels, in particular, benefit from multiple bandwidths because the directional asymmetry signal is present at different scales depending on the functional form of the structural equations.

\subsection{ANM Residual Channels ($+1.50\%$)}

Adding three ANM residual channels to produce the complete 8-channel tensor (v8b: 76.94\%) improves over the 5-channel version (v5m: 75.44\%) by $+1.50\%$. The improvement is concentrated in Collider ($+6.4\%$) and Consequence of X ($+5.0\%$), classes where detecting the directionality of the causal edge is the primary discriminating challenge. The ANM residual $r_{u \to v} = v - \hat{m}(u)$ is approximately independent of $u$ when $u \to v$ is the true causal direction, producing a flat residual curve. Under the reversed direction, the residual retains systematic nonlinear structure. This directional asymmetry -- flat versus structured -- is visible in the 1D residual curve and is what the Conv1D network learns to detect. The multi-bandwidth formulation ensures this directional signal is captured at multiple smoothness scales, making the ANM channel contribution robust across different functional forms.

% =====================================================================
\section{Structural Attention Bias}
\label{sec:struct_results}
% =====================================================================

This section analyzes the single most impactful architectural change: the topology-aware attention mechanism that injects causal structural priors using only six learned scalars.

\subsection{Result and Mechanism ($+2.65\%$)}

Adding structural attention bias to v8b (v11: 79.59\% vs.\ v8b: 76.94\%) yields $+2.65\%$, the largest gain from any single architectural modification across all 30+ experiments conducted in this thesis. The gain is distributed across classes, with the largest improvements in Confounder ($+4.46\%$), Collider ($+3.40\%$), and Mediator ($+1.04\%$). These three classes are precisely the ones whose definitions involve triangle motifs -- chain ($X \to K \to Y$), fork ($K \to X, K \to Y$), and collider ($X \to K \leftarrow Y$) -- that are directly encoded in the six topological relationship types used in the structural attention bias. This distributed gain confirms that the bias is acting on the correct structural feature.

The mechanism operates during attention aggregation: edges that share an endpoint or form a directed chain through a node of interest receive differential attention weights through the learned scalar bias, rather than being treated identically to unrelated background edges. For example, in a Mediator chain $X \to K \to Y$, the edges $X \to K$ and $K \to Y$ form a forward chain relationship (Type 4 in Table~\ref{tab:struct_types}) and receive a shared bias that strengthens their mutual attention, propagating causal context along the pathway.

\subsection{Why Structured Inductive Bias Outperforms Unconstrained Capacity}

Versions v3, v4, v6, v9, and v9b all attempt architectural additions that add unconstrained parameters -- parallel observation cross-attention branches, dual-path transformers, observation-level transformers, node-centric aggregation variants -- and all produce results at or below the 3-channel baseline. The structural attention bias (v11, $+2.65\%$) is the only successful architectural modification among all these experiments.

The distinguishing factor is not parameter count but the alignment of inductive bias with domain knowledge. The six topological relationship types encode explicit prior knowledge about causal motifs (chain, fork, collider, reverse) that the model should attend to. Six additional scalars with this structured prior outperform thousands of additional unconstrained parameters by factors exceeding 3.5 percentage points. This validates the principle that for this problem, structured inductive bias consistently outperforms increased unconstrained capacity. The structured prior does not learn a new representation from scratch; it only learns the relative priority among topologically distinct edge pairs, which is precisely the information encoded in the causal DAG structure.

The lesson extends beyond this specific task: when the domain has known structural constraints (as causal graphs do), encoding those constraints as architectural priors with minimal parameter cost is more effective than adding representational capacity and hoping the model discovers those constraints from data alone.

% =====================================================================
\section{2D Scatter Density Pipeline}
\label{sec:2d_results}
% =====================================================================

This section presents the ablation studies for the 2D scatter density pipeline, establishing that the edge context is structurally necessary, the 1D and 2D representations are complementary, and the final model achieves 81.19\% local and 81.082\% private Balanced Accuracy.

\subsection{Edge Context Is Structurally Necessary}

Table~\ref{tab:edge_context_necessity} presents the ablation evidence for the necessity of edge context in the dual-pipeline architecture. When the entire 1D edge pipeline is removed and nodes are classified from their 2D visual embeddings alone (v26c), Balanced Accuracy collapses to 51.51\% -- a 28.96-point drop from the full model's 80.47\%. A variant using mean pooling of all node embeddings as a substitute for edge context (v28b) achieves only 46.39\%. These results confirm that the 2D density image captures only the local distributional structure of $K$ relative to $X$ and $Y$ but cannot reason about $K$'s causal role within the larger graph structure. The edge encoder, supervised by the binary edge loss against ground-truth adjacency matrices, provides exactly this graph-wide structural context through the four attended edge embeddings ($K \to X$, $K \to Y$, $X \to K$, $Y \to K$). Without it, performance is worse than the ML baseline, confirming that the edge head is a structural necessity rather than auxiliary regularization.

\begin{table}[ht]
	\centering
	\caption{Ablation on the necessity of edge context for node classification. All three variants use the same 8-channel $32\times32$ node images; only the source of graph-wide structural context differs. The full model (v26e 8ch) combines edge context with visual features, achieving 80.47\%. Removing the edge pipeline entirely (v26c) collapses performance by 28.96 points to 51.51\%, worse than the ML baseline. Replacing edge context with mean-pooled node embeddings (v28b) further degrades to 46.39\%. The binary edge loss, which supervises the edge embeddings against ground-truth adjacency entries, is essential for maintaining causally well-structured representations that the node merger depends on.}
	\label{tab:edge_context_necessity}
	\begin{tabular}{llc}
		\toprule
		\textbf{Version}            & \textbf{Graph Context Source}                    & \textbf{Local BA} \\
		\midrule
		v26c                        & None (pure 2D visual, no edge pipeline)        & 51.51\%           \\
		v28b                        & Mean pooling of all node embeddings              & 46.39\%           \\
		v26e 8ch (full model)       & Edge encoder + structural self-attention         & \textbf{80.47\%}  \\
		\bottomrule
	\end{tabular}
\end{table}

\subsection{1D and 2D Representations Are Complementary}

Table~\ref{tab:1d_vs_2d} shows the per-class breakdown comparing the 1D edge-only model (v11: 79.59\%) with the full dual-pipeline model (v26e 8ch: 80.47\%). The 2D pipeline gains most on Confounder ($+2.78\%$), Mediator ($+2.28\%$), and Consequence of Y ($+2.18\%$) -- classes whose distinguishing feature is the joint distributional pattern of $K$ with respect to both $X$ and $Y$ simultaneously. The 1D edge tensor captures each pair separately as marginal curves; the 2D image captures the joint view of $K$ against $X$ and $Y$ together, revealing patterns such as correlated variation across both conditioning variables that are invisible in any single 1D curve. Collider drops slightly ($-2.02\%$), suggesting that the 1D ANM residual channels -- which detect directional asymmetry in marginal pairwise relationships -- are more powerful than 2D density for detecting the collider signature (induced dependence under conditioning). This class-level complementarity confirms that the two representations capture structurally different information about each node's causal role.

\begin{table}[ht]
	\centering
	\caption{Per-class recall comparison between the 1D edge-only model (v11, 79.59\%) and the dual-pipeline model (v26e 8ch, 80.47\%), evaluated on the base dataset without XY augmentation. The 2D pipeline provides the largest gains on Confounder ($+2.78\%$), Mediator ($+2.28\%$), and Consequence of Y ($+2.18\%$), where joint distributional patterns with both $X$ and $Y$ are more informative than marginal pairwise curves. The 1D pipeline outperforms 2D on Collider ($-2.02\%$) because the ANM residual channels capture directional asymmetry more effectively than 2D density for this class, which is distinguished by conditional induction of dependence rather than joint distributional structure. Neither representation fully dominates the other, confirming the complementary nature of 1D pairwise and 2D joint representations.}
	\label{tab:1d_vs_2d}
	\begin{tabular}{lcc}
		\toprule
		\textbf{Class}              & \textbf{v11 (1D only)} & \textbf{v26e 8ch (1D+2D)} \\
		\midrule
		Confounder                 & 75.81\%                & 78.59\%                   \\
		Collider                   & 77.80\%                & 75.78\%                   \\
		Mediator                   & 67.29\%                & 69.57\%                   \\
		Cause of X                 & 82.91\%                & 82.40\%                   \\
		Cause of Y                 & 83.54\%                & 84.71\%                   \\
		Consequence of X           & 79.69\%                & 80.20\%                   \\
		Consequence of Y           & 81.96\%                & 84.14\%                   \\
		Independent                & 87.67\%                & 88.38\%                   \\
		\midrule
		\textbf{Balanced Accuracy} & \textbf{79.59\%}       & \textbf{80.47\%}          \\
		\bottomrule
	\end{tabular}
\end{table}

\subsection{Loss Lambda Ablation}

Table~\ref{tab:lambda_ablation} shows the effect of the edge loss weight $\lambda$ on the final model (v26e 8ch, no XY augmentation). The configuration with $\lambda = 0.7$ (80.47\%) outperforms both $\lambda = 0.3$ (node-favored, 78.77\%) and $\lambda = 1.0$ (edge-only, 80.37\%).

The large gap between $\lambda = 0.3$ and $\lambda = 0.7$ ($-1.70\%$) reflects a training dynamics effect: when the node loss gradient dominates early in training, the edge pipeline receives insufficient supervision and produces inconsistent edge embeddings that degrade the quality of the node merger. The edge pipeline must remain well-supervised throughout training because the node classifier depends entirely on the edge embeddings for graph-wide structural context. The small gap between $\lambda = 1.0$ and $\lambda = 0.7$ ($-0.10\%$) indicates that any configuration with sufficiently high edge weight is adequate; $\lambda = 0.7$ is chosen as it empirically performs best.

\begin{table}[ht]
	\centering
	\caption{Effect of loss weight $\lambda$ on v26e 8ch (base, no XY augmentation). The total loss is $\mathcal{L} = \lambda \cdot \mathcal{L}_{\mathrm{edge}} + (1-\lambda) \cdot \mathcal{L}_{\mathrm{node}}$. Node-favored training ($\lambda=0.3$) produces degenerate edge embeddings that degrade node merger quality, causing a 1.70-point drop. Edge-only training ($\lambda=1.0$) slightly underperforms $\lambda=0.7$ by 0.10 points, suggesting some node gradient is beneficial for stabilizing the fusion layer. The optimal $\lambda=0.7$ reflects the need to keep the edge pipeline well-supervised throughout training, as the node classifier depends entirely on edge embeddings for graph-wide structural context.}
	\label{tab:lambda_ablation}
	\begin{tabular}{lcc}
		\toprule
		\textbf{Configuration} & $\lambda$ & \textbf{Local BA} \\
		\midrule
		v26d node-favored        & 0.3       & 78.77\%           \\
		v26d edge=node           & 1.0       & 80.37\%           \\
		v26e (proposed)          & 0.7       & \textbf{80.47\%}  \\
		\bottomrule
	\end{tabular}
\end{table}

\subsection{Final Model Per-Class Performance}

Table~\ref{tab:perclass_final} presents the per-class recall of the final proposed model (v26e 8ch + XY augmentation) on the local test set. The three hardest classes -- Confounder (74.93\%), Mediator (73.71\%), and Collider (78.70\%) -- remain below the overall average, consistent with competition-wide findings~\cite{olivetti2025adialab}. These are precisely the classes where pairs of roles share similar conditional independence signatures (Confounder vs Cause of X; Mediator vs Consequence of X; Mediator vs Confounder) and where scalar features are most inadequate. The per-class improvements over the reimplementation baseline (v2: 73.96\%) are largest for Mediator ($+13.7\%$), Collider ($+12.7\%$), and Consequence of X ($+8.6\%$), confirming that the 8-channel tensor and structural attention bias specifically address the mechanistic challenges that make these classes difficult.

\begin{table}[ht]
	\centering
	\caption{Per-class recall and sample counts for the final proposed model (v26e 8ch + XY augmentation) on the 1,880-instance local test set. The three hardest classes -- Mediator (73.71\%), Confounder (74.93\%), and Collider (78.70\%) -- are precisely the classes where role pairs share similar conditional independence signatures and require graph-wide structural context to resolve. Independent (88.38\%) is the easiest class because its defining property (no edges to $X$ or $Y$) is detectable from marginal independence alone. The final model achieves 81.19\% Balanced Accuracy locally, 80.96\% on the public leaderboard, and 81.082\% on the private test set.}
	\label{tab:perclass_final}
	\begin{tabular}{lcc}
		\toprule
		\textbf{Class}              & \textbf{Recall}  & \textbf{Count ($n$)} \\
		\midrule
		Confounder                   & 74.93\%          & 682                  \\
		Collider                     & 78.70\%          & 446                  \\
		Mediator                     & 73.71\%          & 483                  \\
		Cause of X                   & 83.32\%          & 983                  \\
		Cause of Y                   & 85.45\%          & 2145                 \\
		Consequence of X             & 81.29\%          & 1571                 \\
		Consequence of Y             & 83.72\%          & 682                  \\
		Independent                  & 88.38\%          & 4494                 \\
		\midrule
		\textbf{Balanced Accuracy}  & \textbf{81.19\%} &                      \\
		\bottomrule
	\end{tabular}
\end{table}

% =====================================================================
\section{Incremental Ablation}
\label{sec:full_ablation}
% =====================================================================

Table~\ref{tab:ablation_full} summarizes the cumulative contribution of each proposed component as incremental gains over the baseline, enabling direct comparison of the relative impact of each contribution.

\begin{table}[ht]
	\centering
	\caption{Incremental ablation: cumulative Balanced Accuracy gains from baseline (v2: 73.96\%) to final model (81.19\%, +7.23 points). Bottom section: 2D pipeline contribution (+0.88\%) measured against 1D-only v11 base, not directly incremental to the v11 row above.}
	\label{tab:ablation_full}
	\begin{tabularx}{\textwidth}{lXcc}
		\toprule
		\textbf{Configuration}      & \textbf{Description}            & \textbf{Local BA} & \textbf{$\Delta$} \\
		\midrule
		v2: baseline                & 3ch reimplementation             & 73.96\%          & ---                \\
		v5m                        & +multi-bw kernel (5ch)           & 75.44\%          & $+1.48$            \\
		v8b                        & +ANM residual ch.\ (8ch)         & 76.94\%          & $+1.50$            \\
		v11                        & +structural attention bias       & 79.59\%          & $+2.65$            \\
		v11+xyaug                  & +XY augmentation                 & 81.14\%          & $+1.55$            \\
		\midrule
		v26e 8ch (1D+2D)           & +2D scatter density             & 80.47\%          & $+0.88^*$          \\
		v26e+xyaug                 & final model                      & \textbf{81.19\%} & $+0.72$            \\
		\bottomrule
	\end{tabularx}
	\medskip
	\noindent *v26e base vs.\ v11 base: 1D+2D vs.\ 1D-only (+0.88\%).
\end{table}

% =====================================================================
\section{Synthesis and Lessons Learned}
\label{sec:lessons}
% =====================================================================

This section extracts four general principles from the experimental results that extend beyond the specific ADIA Lab task to the broader problem of causal discovery from observational data.

\subsection{Shape Information Cannot Be Recovered After Scalar Compression}

The v13 ML fullstack experiment establishes this as the central finding of the thesis. 300+ features, four causal discovery algorithms, and three gradient boosting models together achieve 72.64\% -- lower than every Conv1D-based model that operates on even the minimal 3-channel functional curve representation. The functional form of the causal relationship -- its curvature, multimodality, heteroscedasticity, tail behavior, and directional asymmetry -- contains discriminative information that is not accessible through any scalar summary. Once observations are reduced to aggregate statistics, this information is irrevocably lost. This has direct implications for the design of future causal discovery systems: when the task requires distinguishing roles that differ in the shape of conditional relationships, scalar feature engineering is insufficient regardless of the number of features or the sophistication of the classifier.

\subsection{Structured Inductive Bias Consistently Outperforms Unconstrained Capacity}

The structural attention bias (v11, $+2.65\%$) is the only successful architectural modification among all 30+ experiments conducted in this thesis. Every other architectural addition -- parallel observation branches, dual-path transformers, observation-level transformers, node-centric aggregation, deep sets -- produces results at or below the baseline. The distinguishing factor is not parameter count but the alignment of inductive bias with domain knowledge. Six learnable scalars encoding six topological relationship types outperform thousands of unconstrained parameters by over 3.5 percentage points. This validates the principle that for problems with known structural constraints, encoding those constraints as architectural priors with minimal parameter cost is more effective than adding representational capacity and hoping the model discovers the constraints from data alone. The structured prior does not learn a new representation -- it only learns the relative priority among topologically distinct edge pairs.

\subsection{1D Curves and 2D Density Images Capture Complementary Information}

The 1D edge tensor representation captures directional pairwise relationships between variable pairs as marginal curves, encoding causal asymmetry through ANM residuals. The 2D density image captures the joint distributional pattern of a node with respect to both $X$ and $Y$ simultaneously, encoding joint variation patterns that are invisible in any single 1D marginal curve. Neither representation fully subsumes the other: the 1D pipeline is stronger for Collider (where directional residual asymmetry is the discriminating signal), while the 2D pipeline is stronger for Confounder, Mediator, and Consequence of Y (where joint distributional structure with both $X$ and $Y$ distinguishes the roles). The fusion of both representations (v26e: 80.47\%) outperforms either alone (v11 1D-only: 79.59\%; v26c 2D-only: 51.51\%).

\subsection{Edge Context Is the Structural Foundation, Not an Auxiliary}

Despite the ``auxiliary'' label attached to the binary edge head, the edge pipeline is the structural foundation of the entire model. The 28.96-point collapse of v26c (51.51\%) when the edge pipeline is removed demonstrates that no amount of visual information in the 2D density images can compensate for the absence of graph-wide context. The edge encoder, supervised by the binary edge loss against ground-truth adjacency entries, forces the edge embeddings to encode causally meaningful structure -- specifically, whether each directed edge exists in the ground-truth DAG -- that the node merger requires to determine each node's role. The binary edge loss is not regularization; it is essential supervised context encoding that maintains the discriminative quality of edge embeddings throughout training. Removing this supervision produces degenerate edge representations that degrade the node merger regardless of the quality of the node images.

% =====================================================================
\section*{Chapter Summary}
\addcontentsline{toc}{section}{Chapter Summary}
% =====================================================================

This chapter presented the complete experimental analysis of the proposed method. Six key findings emerged from this work. First, scalar feature compression is irreversible and structurally insufficient for causal role classification: even 300+ features and three ensemble models achieve only 72.64\%, below every Conv1D model. Second, multi-bandwidth kernel regression and ANM residual channels each provide consistent approximately $+1.5\%$ improvements through different mechanisms -- smoothness scale and directional asymmetry respectively. Third, structural attention bias provides the largest single architectural gain of $+2.65\%$ through six learnable scalars encoding topology-aware priors, validating that structured inductive bias outperforms unconstrained capacity by over 3.5 percentage points. Fourth, the 2D scatter density pipeline contributes $+0.88\%$ with complementary class-level improvements, confirming that 1D pairwise and 2D joint representations capture different information. Fifth, edge context is structurally necessary: removing it collapses performance by 29 points, confirming the edge loss is essential supervision rather than auxiliary regularization. Sixth, XY augmentation provides $+0.72\%$ to $+3.32\%$ depending on base model strength, with diminishing returns as the base model becomes richer. The final proposed model achieves 81.19\% locally, 80.96\% on the public leaderboard, and 81.082\% on the private test set, surpassing the competition first-place solution (76.70\%) by $+4.38\%$ and the reimplementation baseline (73.96\%) by $+7.12\%$.